\chapter{Introduction and Project Vision}
% ============================================================

\section{Project Overview}
CARAZONE is a simulation-based federated intrusion detection system for the
Time-Sensitive Networking (TSN) Ethernet backbone of an automotive zonal
electrical/electronic architecture. The system observes TSN protocol-native signals
at each zone controller, extracts timing features over fixed observation windows,
trains a shared detection model across zones using federated learning without
centralising raw feature records, and emits attributed alerts identifying the attack
class and the affected zone.

The delivered artefact is a laboratory validation and research tool, not a production
vehicle component. It comprises a configurable zonal TSN simulation, an attack
injection harness producing ground-truth labels, a feature extraction pipeline, a
federated training and evaluation workflow, and an alert presentation surface.

\section{Problem Statement}
No existing work federates intrusion detection built on TSN's own protocol-native
signals for automotive zonal Ethernet backbones, a topology with heterogeneous
zone-level data that current federated TSN-IDS work (industrial) and automotive IDS
work (CAN bus) both miss.

Concretely: in a zonal architecture the isolation between safety-critical and
non-critical traffic is enforced logically by time-aware scheduling rather than
physically by separate wiring. That scheduling is meaningful only because every node
agrees on the current time. An adversary who corrupts the shared time base therefore
degrades timing guarantees across the backbone without injecting malformed frames,
without raising traffic volume, and without altering any payload. Detection
approaches that observe payload content or flow statistics have no visibility of
this failure mode.

\section{Problem Context and Motivation}
\label{sec:problem-context}
Automotive E/E architecture is migrating from roughly 70--100 distributed ECUs on
segmented buses toward a small number of zone controllers linked by a high-bandwidth
Ethernet backbone. TSN provides the determinism that the displaced buses supplied
natively. Two TSN mechanisms are central: IEEE~802.1AS (gPTP) distributes a common
time base, and IEEE~802.1Qbv (TAS) uses that time base to gate traffic classes into
scheduled windows.

Neither mechanism was standardised with security controls. For gPTP specifically,
cryptographic authentication does not close the full threat surface: an asymmetric
delay attack manipulates the \emph{timing} of authentic messages, corrupting the
peer-delay computation on which offset correction depends. Authentication confirms
that a message is genuine; it cannot establish that the message arrived when it
should have.

The consequence is not immediate catastrophic failure. End-to-end protection
mechanisms at the application layer detect missed updates and degrade to fail-safe
states. The realistic impact is stealthy availability denial: forced degradation,
limp-home transitions and driver warnings at a time of the attacker's choosing. This
is a detection problem rather than a prevention problem, and it is the problem this
project addresses.

\section{Project Vision}
To demonstrate that intrusion detection for automotive zonal TSN backbones is more
effective when built on protocol-native timing signals than on generic traffic
statistics, and that such detection can be trained across heterogeneous zones using
federated learning without centralising raw per-zone observations.

\section{Project Objectives}
\begin{enumerate}
  \item \textbf{O1.} Implement a configurable multi-zone automotive TSN topology in
        OMNeT++/INET with gPTP synchronisation and reproducible seeded execution.
  \item \textbf{O2.} Implement an attack injection harness for TSN timing attacks that
        produces ground-truth labels with zone, time-window and attack-class attribution.
  \item \textbf{O3.} Implement a passive feature extraction pipeline converting raw
        simulation output into labelled, windowed, protocol-native feature vectors.
  \item \textbf{O4.} Train and evaluate a detection model on those features, and
        demonstrate by controlled comparison whether protocol-native features
        outperform generic traffic features on the same scenarios.
  \item \textbf{O5.} Implement federated training across per-zone clients and compare
        it against centralized-pooled and local-only training baselines.
  \item \textbf{O6.} Produce attributed alerts carrying attack class, affected zone,
        time window and confidence, and present them through an operator-facing surface.
\end{enumerate}

Objectives O1--O3 are complete at the mid-report stage. O4--O6 are scheduled for
Iterations~03 and~04.

\section{Project Scope}
\label{sec:scope}

\subsection*{Scope statement}
The project targets the TSN Ethernet backbone connecting zone controllers to central
compute in automotive zonal E/E architectures, focusing specifically on attacks
against TSN features. The project trains a detection model across simulated
zone/platform clients standing in for heterogeneous vehicle architectures (zonal
topologies), using federated learning. Evaluation compares this federated approach
against a centralized version of the same detector and against existing
non-federated detection work, using synthetic attack datasets generated via a TSN
simulation framework.

\subsection*{Included and excluded functionality}

\begin{longtable}{p{0.44\textwidth} p{0.48\textwidth}}
\toprule
\textbf{Included (committed)} & \textbf{Excluded} \\
\midrule
\endhead
Multi-zone TSN topology, simulated in OMNeT++/INET with gPTP synchronisation &
Hardware testbed or physical vehicle validation \\
\addlinespace
Attacks against gPTP synchronisation (delay injection, sync starvation) and against
FRER (IEEE 802.1CB) stream redundancy &
Link-layer attacks outside the TSN control and redundancy planes \\
\addlinespace
Passive extraction of protocol-native timing features over fixed windows &
Active probing, packet modification, or in-path interference \\
\addlinespace
Supervised multiclass detection (benign $+$ attack classes) &
Unsupervised/one-class detection as the primary mechanism \\
\addlinespace
Federated training (FedAvg) with centralized and local-only baselines &
Secure aggregation, differential privacy, model-poisoning defence \\
\addlinespace
Zone-level attack attribution &
ECU-level or port-level localisation \\
\addlinespace
Operator-facing alert presentation surface &
Production vehicle integration or in-cabin driver interface \\
\bottomrule
\end{longtable}

\subsection*{Scope items in progress, and deferred scope}
\textbf{FRER-based detection (in scope, implementation in progress).} Frame
Replication and Elimination for Reliability is committed scope and is under active
development. It is not yet exercised in the dataset because the
four-zone tree topology used for Iteration~1 gives each zone a single path to the
backbone, and FRER requires disjoint redundant paths between talker and listener.
That blocking dependency has since been removed. A five-zone ring reference topology
has been designed, generated and structurally validated (Section~\ref{sec:ref-topology}); its automated
check confirms that cutting any single ring or uplink link partitions nothing, which
is the machine-checkable form of the disjoint-path requirement. The remaining work is
therefore bounded and specified: configuring the 802.1CB sequence generation and
recovery functions at the talker and listener ends of the redundant streams;
injecting stream-redundancy manipulation as an attack; and adding the FRER feature
family to the extraction pipeline: inter-member arrival skew, sequence-number
deviation between redundant members, and inter-arrival cadence per member. The
attack model is stream-redundancy manipulation: selective suppression or delay of one
member of a replicated stream, which degrades the reliability guarantee while every
frame remains individually authentic. This family is attractive for the project's
central claim because its features are structurally independent of the gPTP offset
residuals, so it tests whether the cross-zone residual principle generalises beyond
the synchronisation plane. Scheduling and the exit criterion are given in
Section~\ref{sec:remaining-iterations}. The topology dependency, previously the single blocker, is resolved in
design and awaits a baseline simulation run.

\textbf{Schedule-window (802.1Qbv) violation detection (deferred).} TAS gate scheduling is not
yet configured in the topology, so gate-conformance features cannot be computed.
This is retained as a planned extension because it supplies a second, independent
protocol-native feature family and a direct physical consequence measure (deadline
misses).

The 802.1Qbv deferral follows the scope-protection rule established in the project
proposal, which identifies scope expansion through additional attack classes as a
high-impact risk and names extension features as the first items to drop. FRER is
distinguished from that rule because it is a committed objective with a bounded
dependency rather than an opportunistic addition.

\section{Project Modules}

\begin{longtable}{p{0.05\textwidth} p{0.26\textwidth} p{0.57\textwidth}}
\toprule
\textbf{\#} & \textbf{Module} & \textbf{Responsibility} \\
\midrule
\endhead
M1 & Zonal Topology Simulation & Constructs the multi-zone TSN network model:
grandmaster, backbone switch, zone controllers, endpoint ECUs, central vehicle
compute; gPTP synchronisation; seeded reproducible execution. \\
\addlinespace
M2 & Attack Injection \& Dataset & Injects TSN protocol-native attacks with
configurable start time, intensity and target; emits ground-truth labels with zone,
window and class attribution. \\
\addlinespace
M3 & Feature Extraction Layer & Passively converts raw per-node simulation output
into fixed-size windowed feature vectors using protocol-native signals rather than
generic traffic statistics. \\
\addlinespace
M4 & Federated Detection Model & Trains a shared detection model across per-zone
clients; produces the global model distributed back to zones for inference. \\
\addlinespace
M5 & Zone Participation \& Aggregation & Governs how zones with differing traffic
profiles contribute to the shared federated model. \\
\addlinespace
M6 & Alerting \& Presentation & Packages classifications into attributed alerts and
presents them through an operator-facing incident queue. \\
\bottomrule
\end{longtable}

Modules M1--M3 are implemented. M4--M6 are designed and scheduled.

\section{Stakeholders and User Classes}
\label{sec:stakeholders}

\begin{longtable}{p{0.24\textwidth} p{0.34\textwidth} p{0.34\textwidth}}
\toprule
\textbf{Stakeholder} & \textbf{Need} & \textbf{Evidence required} \\
\midrule
\endhead
Automotive network security / validation engineer &
Determine whether timing attacks are detectable in a candidate zonal TSN
configuration &
Attributed alerts with class, zone, time window, confidence; reproducible scenarios \\
\addlinespace
Vehicle Security Operations Centre (vSOC) analyst &
Triage incoming incidents across a fleet and distinguish fault from campaign &
Incident queue with severity, zone, attack class, confidence and status \\
\addlinespace
TSN validation engineer &
Confirm that monitoring does not disturb protected traffic &
Matched monitoring-enabled vs.\ monitoring-disabled timing measurements \\
\addlinespace
Project supervisor / evaluator &
Judge whether the approach is technically justified &
Baseline comparison, reproducible experiments, documented failure analysis \\
\addlinespace
Future researcher &
Reproduce or extend the work &
Versioned configuration, seeds, raw outputs, run instructions \\
\bottomrule
\end{longtable}

\textbf{Operating context.} The intended deployment position of the detector is a
zone controller within the vehicle, with aggregation on the central vehicle compute
unit. In this project the entire system is exercised in simulation on a development
workstation; no claim of in-vehicle deployment is made.

\section{Success Criteria}
\label{sec:success-criteria}

Success is assessed against measurable criteria rather than a single aggregate
accuracy figure. Class imbalance in the generated data (see Section~\ref{sec:test-results}) makes
accuracy uninformative, so it is excluded as a headline measure.

\begin{longtable}{p{0.06\textwidth} p{0.34\textwidth} p{0.52\textwidth}}
\toprule
\textbf{ID} & \textbf{Criterion} & \textbf{Measure and target} \\
\midrule
\endhead
SC-1 & Attack effects are measurable and repeatable &
Injected attacks produce a timing signature distinguishable from benign behaviour
across $\geq$3 seeds, with ground-truth labels. \textbf{Met} for delay injection. \\
\addlinespace
SC-2 & Cross-zone features outperform within-zone features &
Identical model trained on (a) cross-zone residual and (b) within-zone statistics over
identical run-level splits; macro-F1 of (a) exceeds (b) by a margin that holds across
seeds and across models. \textbf{Met}: $+0.43$ to $+0.49$ macro-F1 over three models
(Section~\ref{sec:acceptance-validation}). \\
\addlinespace
SC-3 & Detection is operationally timely &
Detection latency, reported in synchronisation intervals from attack onset.
\textbf{Partially met}: two intervals under a three-window confirmation rule
(Section~\ref{sec:acceptance-validation}). Comparison against a schedule margin awaits time-aware shaping. \\
\addlinespace
SC-4 & Federated training is justified &
FedAvg macro-F1 exceeds local-only training, and the gap to centralized-pooled
training is quantified. A null result is reported as a finding, not concealed. \\
\addlinespace
SC-5 & Attribution is correct &
Zone attribution accuracy (top-1) reported separately from detection performance. \\
\addlinespace
SC-6 & Alert volume is operationally acceptable &
False alarms per vehicle-hour reported at the chosen operating point, in addition to
false-positive rate. \\
\addlinespace
SC-7 & Results are reproducible &
A documented configuration and seed regenerate the same scenario structure and
ground-truth labels on a clean environment. \\
\bottomrule
\end{longtable}

\section{Constraints and Assumptions}
\label{sec:assumptions}

\subsection*{Constraints}
\begin{itemize}
  \item \textbf{C1.} Evaluation is simulation-based. No hardware TSN testbed or
        vehicle platform is available, so all traffic and all attacks are synthetic.
  \item \textbf{C2.} Observable signals are limited to those the simulation framework
        exposes. Features depending on internal servo state must be instrumented
        explicitly, and availability is verified per feature before commitment.
  \item \textbf{C3.} Development and simulation run on a single workstation; the
        federated setting is simulated (clients as processes), not physically
        distributed.
  \item \textbf{C4.} Two-member team over two semesters, which bounds the number of
        attack classes and model variants that can be implemented and defended.
\end{itemize}

\subsection*{Assumptions}
\begin{itemize}
  \item \textbf{A1. Static network environment.} The topology, the gPTP domain
        configuration and (when introduced) the gate schedule are statically
        configured and do not reconfigure at runtime. This assumption is load-bearing:
        the detection approach uses deviation from a known, configured expectation,
        which is available \emph{a priori} only in a static environment. In a
        dynamically reconfigured network, legitimate reconfiguration events would
        present as anomalies and would require explicit handling. Dynamic
        reconfiguration is out of scope and recorded as a limitation.
  \item \textbf{A2. Static (one-shot) training.} Federated training occurs once,
        offline, before deployment. After the global model is distributed, clients
        perform inference only; no online retraining occurs. Continuous federated
        learning would require solving live ground-truth labelling in deployment,
        which is out of scope.
  \item \textbf{A3.} The grandmaster is statically assigned and trusted. Attacks
        originate from a compromised on-path node within a zone, not from the time
        source itself. Grandmaster-level compromise is addressed as a scope boundary
        in Section~\ref{sec:technical-challenges}.
  \item \textbf{A4.} The aggregation server (central vehicle compute) is trusted.
        Secure aggregation and model-poisoning resistance are out of scope.
  \item \textbf{A5.} Monitoring is passive and sits outside the forwarding path.
\end{itemize}

\section{Ethical, Legal and DEI Considerations}

\textbf{Ethical.} The project uses only synthetic, simulated data. No real vehicle
data, no OEM data and no personal data are collected or processed. No human subjects
are involved. Attack implementations target a simulated network only; no capability
is developed against any real or deployed system.

\textbf{Legal.} No proprietary OEM communication matrices or topology data are used.
OMNeT++ and INET are used under their respective open-source licences, which are
recorded in the dependency register (Section~\ref{sec:dependencies}). Published datasets, if
incorporated for validation, will be used under their stated academic-use terms with
required citation.

\textbf{DEI and bias.} Synthetic topologies are constructed from published
architectural patterns rather than any specific manufacturer's proprietary design,
which avoids embedding a single vendor's assumptions into the evaluation. The
project produces no output that classifies or affects individuals.

\textbf{Dual-use consideration.} The attack implementations are confined to the
simulation environment and are described at the level of effect rather than as a
reusable exploit toolkit.

% ============================================================